\documentclass[14pt,a4paper,oneside]{extbook}
\usepackage[utf8]{vietnam}
\usepackage{amsmath,amssymb,mathrsfs,fancyhdr,enumerate,multirow,makecell}
\usepackage{tkz-euclide,tikz-3dplot,tikz,tkz-tab}
\usetikzlibrary{arrows.meta,calc,intersections,angles,quotes,backgrounds,shapes.geometric,patterns,through,shadows,positioning}
\usepackage{yhmath}
\usepackage[top=0.7cm, bottom=0.7cm, left=1.0cm, right=1.0cm, headheight=16pt, headsep=3mm, footskip=5mm]{geometry}
\usepackage[hidelinks,unicode]{hyperref}
\usepackage{qrcode}
\usepackage[many]{tcolorbox}
\usepackage{fontawesome}

\pagestyle{fancy}
\fancyhf{}
\fancyhead[L]{\footnotesize\bfseries\color{teal!85!black}\faBook~SỔ TAY KIẾN THỨC TOÁN 4}
\fancyhead[R]{\footnotesize\bfseries\color{blue!85!black}Lớp Toán Thầy Lĩnh -- loptoanthaylinh.vn}
\fancyfoot[C]{\footnotesize\bfseries\color{gray!80!black}-- Sổ tay kiến thức: Bài 12. Các số trong phạm vi lớp triệu --}
\renewcommand{\headrulewidth}{0.4pt}

\begin{document}

% --- BANNER TIÊU ĐỀ BÀI HỌC ---
\noindent
\begin{tcolorbox}[
  enhanced,
  colback=teal!8!white,
  colframe=teal!80!black,
  arc=3mm, boxrule=1pt,
  left=3mm, right=3mm, top=2mm, bottom=2mm,
  halign=center
]
  {\Large\bfseries\color{teal!90!black}\faBookmark\ BÀI 12: CÁC SỐ TRONG PHẠM VI LỚP TRIỆU}\\[1.5pt]
  {\footnotesize\itshape\color{black!75} Hàng triệu, chục triệu, trăm triệu, đọc viết số có 7, 8, 9 chữ số và cấu tạo số}
\end{tcolorbox}

\vspace{2.5mm}

% --- KHỐI 1 & 2: LÝ THUYẾT & KĨ NĂNG ---
\noindent
\begin{minipage}[t]{0.49\textwidth}
  \begin{tcolorbox}[
    enhanced,
    title={\footnotesize\bfseries\color{white}\faGlobe\ 1. Lớp triệu \& Các hàng},
    colback=blue!2!white,
    colframe=blue!75!black,
    arc=2.5mm, boxrule=0.8pt, top=1.8mm, bottom=1.8mm, left=2mm, right=2mm
  ]
    \footnotesize
    $\bullet$ \textbf{Lớp triệu gồm 3 hàng:}\\
    - Hàng triệu ($1\,000\,000$ -- 7 chữ số).\\
    - Hàng chục triệu ($10\,000\,000$ -- 8 chữ số).\\
    - Hàng trăm triệu ($100\,000\,000$ -- 9 chữ số).\\
    $\bullet$ \textbf{Quan hệ giữa các hàng liền kề:}\\
    $10\text{ triệu} = 1\text{ chục triệu}$; $10\text{ chục triệu} = 1\text{ trăm triệu}$.\\
    $\bullet$ Số nhỏ nhất có 7 chữ số: $1\,000\,000$.
  \end{tcolorbox}
\end{minipage}%
\hfill
\begin{minipage}[t]{0.49\textwidth}
  \begin{tcolorbox}[
    enhanced,
    title={\footnotesize\bfseries\color{white}\faPencil\ 2. Kĩ năng đọc viết số lớp triệu},
    colback=teal!2!white,
    colframe=teal!75!black,
    arc=2.5mm, boxrule=0.8pt, top=1.8mm, bottom=1.8mm, left=2mm, right=2mm
  ]
    \footnotesize
    $\bullet$ \textbf{Tách lớp từ phải sang trái:}\\
    (Lớp triệu) (Lớp nghìn) (Lớp đơn vị).\\
    $\bullet$ \textbf{Cách đọc:} Đọc từng lớp từ trái sang phải kèm theo tên lớp.\\
    \textit{Ví dụ:} $315\,240\,806$ đọc là: Ba trăm mười lăm \textbf{triệu} hai trăm bốn mươi \textbf{nghìn} tám trăm linh sáu.\\
    $\bullet$ Khi viết, để khoảng cách nhỏ giữa các lớp.
  \end{tcolorbox}
\end{minipage}

\vspace{2.5mm}

% --- KHỐI 3: SƠ ĐỒ TƯ DUY ---
\noindent
\begin{tcolorbox}[
  enhanced,
  title={\footnotesize\bfseries\color{white}\faSitemap\ Sơ đồ tư duy: Cấu trúc 3 lớp của số trong phạm vi lớp triệu},
  colback=yellow!2!white,
  colframe=teal!85!black,
  arc=3mm, boxrule=0.9pt, top=1.8mm, bottom=1.8mm, left=1.5mm, right=1.5mm,
  halign=center
]
  \begin{tikzpicture}[
    scale=0.82,
    every node/.style={transform shape},
    root/.style={rectangle, rounded corners=5pt, fill=teal!85!black, text=white, font=\bfseries\footnotesize, inner sep=4.5pt, align=center},
    cardT/.style={rectangle, rounded corners=4pt, fill=purple!12, draw=purple!80!black, text=purple!95!black, font=\scriptsize, inner sep=3.5pt, align=center, minimum width=3.8cm},
    cardN/.style={rectangle, rounded corners=4pt, fill=blue!10, draw=blue!75!black, text=blue!90!black, font=\scriptsize, inner sep=3.5pt, align=center, minimum width=3.8cm},
    cardD/.style={rectangle, rounded corners=4pt, fill=teal!12, draw=teal!80!black, text=teal!95!black, font=\scriptsize, inner sep=3.5pt, align=center, minimum width=3.8cm}
  ]
    \node[root] (core) at (0, 1.15) {\faDatabase\ SỐ TRONG PHẠM VI LỚP TRIỆU (ĐẾN 9 CHỮ SỐ)};

    \node[cardT] (c1) at (-4.8, -0.35) {\textbf{LỚP TRIỆU}\\[1.5pt]$\bullet$ Hàng trăm triệu\\$\bullet$ Hàng chục triệu\\$\bullet$ Hàng triệu};
    \node[cardN] (c2) at (0, -0.35) {\textbf{LỚP NGHÌN}\\[1.5pt]$\bullet$ Hàng trăm nghìn\\$\bullet$ Hàng chục nghìn\\$\bullet$ Hàng nghìn};
    \node[cardD] (c3) at (4.8, -0.35) {\textbf{LỚP ĐƠN VỊ}\\[1.5pt]$\bullet$ Hàng trăm\\$\bullet$ Hàng chục\\$\bullet$ Hàng đơn vị};

    \draw[teal!85!black, thick, -{Stealth[length=2mm]}] (core.south) -- ++(0, -0.2) -| (c1.north);
    \draw[teal!85!black, thick, -{Stealth[length=2mm]}] (core.south) -- (c2.north);
    \draw[teal!85!black, thick, -{Stealth[length=2mm]}] (core.south) -- ++(0, -0.2) -| (c3.north);
  \end{tikzpicture}
\end{tcolorbox}

\vspace{2.5mm}

% --- KHỐI 4: MẸO & BẪY ---
\noindent
\begin{minipage}[t]{0.54\textwidth}
  \begin{tcolorbox}[
    enhanced,
    title={\footnotesize\bfseries\color{white}\faKey\ Phân tích số thành tổng các triệu, nghìn},
    colback=green!3!white,
    colframe=green!60!black,
    arc=2.5mm, boxrule=0.8pt, top=1.8mm, bottom=1.8mm, left=2.2mm, right=2.2mm
  ]
    \footnotesize
    $\bullet$ \textbf{Ví dụ phân tích:}\\
    $5\,240\,100 = 5\,000\,000 + 200\,000 + 40\,000 + 100$\\
    $\bullet$ Có thể viết gộp theo từng lớp:\\
    $5\,240\,100 = 5\text{ triệu} + 240\text{ nghìn} + 100\text{ đơn vị}$.\\
    $\bullet$ Chữ số 0 ở hàng nào thì không cần cộng giá trị hàng đó vào tổng.
  \end{tcolorbox}
\end{minipage}%
\hfill
\begin{minipage}[t]{0.44\textwidth}
  \begin{tcolorbox}[
    enhanced,
    title={\footnotesize\bfseries\color{white}\faExclamationTriangle\ Mẹo tránh bẫy trắc nghiệm},
    colback=red!2!white,
    colframe=red!70!black,
    arc=2.5mm, boxrule=0.8pt, top=1.8mm, bottom=1.8mm, left=2.2mm, right=2.2mm
  ]
    \footnotesize
    $\bullet$ \textbf{Bẫy cả lớp bằng 0:}\\
    Ví dụ số $7\,000\,015$: Lớp nghìn toàn số 0 thì \textbf{không đọc chữ nghìn}, chỉ đọc: "Bảy triệu không trăm mười lăm".\\
    $\bullet$ \textbf{Bẫy viết thiếu số 0:}\\
    Nghe đọc số phải đếm đủ số chữ số của từng lớp (mỗi lớp 3 chữ số).
  \end{tcolorbox}
\end{minipage}

\end{document}
